\documentclass[10pt,a4paper]{amsart}
\usepackage[margin=19mm]{geometry}
\usepackage{amsmath,amssymb,mathtools}
\usepackage[scr=rsfs,cal=euler]{mathalfa}
\usepackage{xfrac}
\usepackage[unicode,hidelinks,bookmarks=false]{hyperref}
\newcommand{\ie}{i.e.\ }
\newcommand{\cf}{cf.\ }
\newcommand{\resp}{respectively\ }
\newcommand{\Subsection}{\subsection}
\providecommand{\nobreakdash}{\nobreak\hbox{-}}
\setlength{\emergencystretch}{2em}
\multlinegap0pt
\usepackage{xcolor}
\usepackage{amssymb}
\usepackage[shortlabels]{enumitem}
\newtheorem{lemma}{Lemma}
\title{$W^*$-categories are von Neumann}
\author{Matthew Daws}
\date{Source: arXiv:2609.30015v1 (CC BY 4.0)}
\begin{document}
\maketitle

\noindent\emph{Source and licence.} $W^*$-categories are von Neumann. Version arXiv:2609.30015v1 (CC BY 4.0), distributed by arXiv under the Creative Commons Attribution 4.0 International licence, http://creativecommons.org/licenses/by/4.0/, as recorded in the arXiv licence metadata for this exact version and shown on the abstract page https://arxiv.org/abs/2609.30015v1. Full licence notice: \texttt{licenses/arxiv-2609.30015-CC-BY-4.0.txt}.\par\smallskip

\noindent\textit{Editorial scope.} Counted unit: the article's lemma lemma (source0.tex lines 161--163). The article's complete original proof of it is delivered with it (source0.tex lines 164--170, one \texttt{proof} environment of 7 lines). No mathematics is written, repaired or re-derived: the statement and the proof are the article's own lines, in the article's own notation, and no retained line is substituted or re-flowed.  No further source-local context is needed: every cross-reference of every retained line resolves inside the two delivered spans.  Nothing was omitted from any retained span, so the package carries no omission record.  No retained line contains a citation, so the wrapper carries no bibliography and no bibliography entry.  No retained line contains a figure environment, an image-inclusion command or drawing code, so no image is delivered and the package declares no asset dependency.  Editorial formatting only, in addition: an AMS \texttt{amsart} preamble that restates the article's own declarations and macros used by the retained lines, namely the environments \texttt{lemma} on separate counters, together with the standard packages those lines need.  The retained environments print in this document as Lemma 1 in order of appearance, whereas the article prints the same items with the numbering of its own full document.  The complete native source of the article as distributed in the arXiv e-print for this exact version is bundled unchanged as \texttt{sources/211717/source0.tex}.
\begin{lemma}
The map $\theta \colon V_{(r)} \to rr^* C rr^*; x \mapsto xr^*$ is a $*$-isomorphism with inverse $a\mapsto ar$.  This shows in particular that $V_{(r)}$ is a $C^*$-algebra.
\end{lemma}

\begin{proof}
As $rr^*$ is a projection in $C$, we see that $rr^* C rr^*$ is a $C^*$-algebra, with unit $rr^*$.  For any $x\in V$, indeed $xr^* \in C$, and when $x\in V_{(r)}$ we have that $rr^*(xr^*)rr^* = rr^* x r^* = xr^*$ so $\theta$ does map $V_{(r)}$ to $rr^* C rr^*$.  Also $\theta(x^\flat) = rx^*rr^* = r(rr^*x)^* = rx^* = \theta(x)^*$ and for $x,y\in V_{(r)}$ we have $\theta(x\bullet y) = xr^*yr^* = \theta(x)\theta(y)$, so $\theta$ is a $*$-homomorphism.

Let $a = yz^*$ for some $y,z\in V$ and consider $x = (rr^* a rr^*) r$.  Then $x = rr^* yz^* r = rr^* (yz^* r) \in V$ by the TRO property.  As $V$ is norm closed, it follows that for any $a\in C$ we have that $x = (rr^* a rr^*) r \in V$, and so for any $a\in rr^* C rr^*$ we have that $x = ar \in V$.  For such an $x$ we now see that $rr^* x = rr^* a r = ar$ and $x r^*r = arr^*r = ar = x$ so that $x\in V_{(r)}$.  Finally, $\theta(x) = arr^* = a$ and so we conclude that $\theta$ is a bijection.

As $\|r^*\|\leq 1$ we see that $\theta$ is a contraction, and similarly $\theta^{-1}$ is a contraction, so $\theta$ is an isometry.  Hence the given norm on $V_{(r)}$ satisfies the $C^*$-condition.
\end{proof}

\end{document}
